\section{Introduction}\label{sec:intro}
An LLM Agent receives a support ticket, checks an order and submits a refund through a third-party commerce service. The traffic is encrypted, yet an observer of the path to that service sees which endpoints are contacted, how large each exchange is and how long the Agent waits between them. Recent work shows that such metadata reveal prompts, tasks, workflows and user traits~\cite{jeong2026wra,dangol2026workflow,pouryousef2026traffic,li2026ant}.

A defender who wants to hide the refund can add decoy calls. Classical traffic-analysis defenses discard their cover at a cooperating endpoint before it reaches an application~\cite{apthorpe2019stp}. On the hop to a third-party service, the defender can still pad TLS records and HTTP/2 frames, which the service's protocol stack discards~\cite{rfc8446,rfc9113}. Padding shapes the defender's requests, but it cannot make the service answer. If the observer may learn which request--response exchanges occur and how long the service takes, decoys must reach the application and elicit real responses. Such \emph{executed cover} has two outputs: a communication trace and application work, which can consume quota, touch other objects and re-enter the planner's context.

\paragraph{The central observation}
Executed cover leaks through the process that sends it. An honest tool call is issued after the model decides on it, so it follows a model wait. A decoy inserted by the controller does not have to wait for anything. We show that this difference is exploitable on real TLS records, that hiding it by sampling historical gaps leaves the signal largely intact, and that the natural alternative, a second planner that generates cover by running a complete independent task, changes the trace at a measurable cost. Decoy-query obfuscation in web search taught a related lesson, that decoys are only as good as their generator~\cite{howe2008trackmenot,peddinti2010trackmenot,balsa2012obpws}. In Agents the generator is coupled to model inference, and decoys perform business work.

\begin{figure*}[t]
\centering
\begin{tikzpicture}[x=1cm,y=1cm,font=\small,>=Latex]
\tikzset{box/.style={draw,rounded corners,align=center,minimum height=.85cm,text width=2.4cm},
 note/.style={align=center,font=\footnotesize,text width=3.1cm}}
\node[box,fill=blue!5] (u) at (1.3,1.15) {User\\``Process the next ticket''};
\node[box,fill=blue!5] (a) at (5.1,1.15) {Agent\\plan, check, act, retry};
\node[box,fill=green!6] (m) at (9.2,2.6) {Model service\\next action / streamed output};
\node[box,fill=green!6] (t) at (9.2,-.35) {Tool service\\ticket, order, refund};
\node[box,fill=orange!8] (d) at (14.1,1.15) {Specialist Agent\\optional delegation};
\draw[<->,thick] (u)--node[above,yshift=15pt]{client link C}(a);
\draw[<->,thick] (a.north east)--node[above,sloped]{model link M}(m.west);
\draw[<->,thick] (a.south east)--node[below,sloped]{tool link T}(t.west);
\draw[<->,thick] (a.east)--node[above]{coordination A}(d.west);
\node[note] at (2.0,-.35) {\textbf{Transit observer}\\lengths, gaps, connections;\\destinations when available};
\node[note] at (13.7,-.35) {\textbf{Recipient / active service}\\received arguments;\\optionally selected responses};
\node[draw=black!45,fill=black!3,align=left,text width=15.7cm,inner sep=6pt] at (7.9,-1.65) {
\textbf{Joint protection decision:} Which observed traces remain distinguishable? Does cover preserve the requested goal, allowed business state and resource budget? State where the cover stops and whether its feedback reaches the planner.};
\end{tikzpicture}
\caption{Running support-ticket example and observation rights. A tool result returns to the Agent and can enter the next model request. Each encrypted link exposes its own traffic shape; a recipient additionally sees the request addressed to it. Concurrent tickets may share a connection on either link. The arrows represent application communication, not a one-to-one mapping from tool events to packets.}
\label{fig:chain}
\end{figure*}


\paragraph{Approach}
We state when cover must reach the application, classify executed cover along four axes (termination, state, feedback and dispatcher), and list the service capabilities each policy assumes. We define an evaluation that binds residual inference, requested-goal completion, business-state integrity and resource cost to the same protected executions. We apply it to 3,000 live executions of two fixed-weight planners, Qwen3:8B and Llama3.1:8B, over MCP and TLS~1.3/HTTP/2 on one host, comparing four cover dispatchers against no cover with a cooperating tool provider.

\paragraph{Findings}
\begin{itemize}
\item \textbf{Controller dispatch is exploitable.} Given call boundaries, label-free gap removal raises recognition of the requested action from 75 to 112 (Qwen) and from 80 to 108 (Llama) of 120 tasks; both gains survive Holm correction, and refitting on independent packet-capture clocks reproduces them.
\item \textbf{History matching hides the shortcut, not the signal.} Dispatching after sampled historical gaps defeats the removal, yet adapted attacks still gain over size-only attacks in three of four comparisons after Holm correction, recognizing 93--106/120.
\item \textbf{Without call boundaries the picture changes.} The prespecified removal attack loses its advantage on complete encrypted connections, while recognition stays well above chance.
\item \textbf{A shadow planner lowers recognition at a cost.} The strongest tested attack reaches 70/120 and 64/120, consistent with a 7/12 plausible-deniability reference; the shadow doubles model tokens and adds more and different cover calls, so the gain is not yet attributable to planner generation alone.
\item \textbf{Goal completion is not state integrity.} In a stress control where cover writes to live state, 95 of 102 completed Qwen tasks leave collateral damage that a goal check misses.
\end{itemize}

\paragraph{Contributions}
(1) The dispatcher as a leakage source for executed cover, with live evidence on TLS records (Sections~\ref{sec:setup}--\ref{sec:results}). (2) A threat model that states when cover must reach the application, a four-axis design space and the service contract each policy assumes (Sections~\ref{sec:threat}--\ref{sec:design}). (3) A joint evaluation protocol with interval-based decisions and the plausible-deniability reference for planner cover (Section~\ref{sec:method}). Appendix~\ref{app:map} maps all forty-four existing Agent side-channel settings; in the four query/workflow privacy settings among them, none reports both an adapted attacker and original-task utility under the same intervention.
